\documentclass[11pt,a4paper]{article}
\usepackage[margin=24mm]{geometry}
\usepackage{amsmath,amssymb,amsthm,booktabs,microtype}
\usepackage[T1]{fontenc}
\usepackage{lmodern}
\usepackage[colorlinks=true,linkcolor=blue,urlcolor=blue,citecolor=blue]{hyperref}
\hypersetup{pdftitle={Subset Sums: a checked continuation from 100 to 99},pdfauthor={Sum-of-Subset research project}}
\newcommand{\FS}{\operatorname{FS}}
\newcommand{\Z}{\mathbb Z}
\newcommand{\T}{\mathcal T}
\newtheorem{theorem}{Theorem}[section]
\newtheorem{lemma}[theorem]{Lemma}
\newtheorem{proposition}[theorem]{Proposition}
\title{Subset Sums in a Fourfold Sumset\\\large A checked continuation from 100 to 99}
\author{Sum-of-Subset research project}
\date{16 October 2026}
\begin{document}
\maketitle
\begin{abstract}
We exclude cardinality 100, continuing the independently checked bound
$N\le100$ to $N\le99$. The new cases allow two empty covering units, or
one empty unit together with two centers outside their target intervals.
Equality in the defect budget supplies a short rigidity lemma for the
second case. The first case uses full translates and weak outer order,
without assuming outer integrality. All 59,206 linear systems are
certified by 204 freshly generated nonnegative integer row combinations.
The verifier checks every coefficient and every integer parameter using
only the Python standard library. The resulting checked range is
$71\le N\le99$; the stronger archived claim $N\le84$ remains outside
the complete independent audit.
\end{abstract}
\tableofcontents

\section{Statement and precise scope}
For a finite real set $A$, define
\[
 \FS(A)=\left\{\sum_{a\in X}a:X\subseteq A\right\}.
\]
The empty sum is included; $4A=A+A+A+A$ permits repeated summands.
Let $N$ be the largest cardinality of a set satisfying $\FS(A)\subseteq4A$.
Write $T_r=r(r+1)/2$.
\begin{theorem}\label{thm:main}
Every feasible finite real set has at most 99 elements.
\end{theorem}
The known feasible 71-element set
\[
 [-31,31]_{\Z}\cup\{\pm93,\pm155,\pm217,\pm713\}
\]
is unchanged, so the independently closed interval becomes $71\le N\le99$.

We use the previously audited theorem $N\le100$ and the structural
lemmas in the 115-to-103 and 103-to-100 reports \cite{base103,base100}.
The size-100 lattice gates are checked again below. The new equality
lemma, coverage of the extra cases, and all exact linear systems are
given explicitly. The archived upper-86 report \cite{archive} contains
the strategy for this endpoint. Our contribution here is a fresh,
compactly specified continuation with newly generated certificates and
a complete current run, rather than a claim to improve on the archived
numerical record 84. No archived endpoint certificate is read by the
new verifier.

\section{Lattice reduction and the size-100 budget}
A hypothetical feasible 100-element set contains zero: otherwise
adjoining zero preserves its subset-sum support and enlarges $4A$, giving
a forbidden 101-element set. Integralize by fixing a four-term
representation for each subset sum, approximating within the kernel of
the resulting homogeneous rational linear equations, preserving strict
coordinate order, and clearing denominators.

Remove the four largest magnitudes of each sign to obtain $D$, then make
two four-extreme peels $D\supset D_1\supset E$ preserving zero. The
inherited cover and peel comparisons give $|\FS(E)|\le24|D|+77$.
The one-sided alternative is excluded by $1+T_{84}>24\cdot100+77$.
If $E$ has $p$ opposite pairs and $q$ unpaired absolute values, the
inherited support estimate gives
\[
 q^2\le96e+1073-e^2,\qquad q\equiv e-1\pmod2,\qquad e=84.
\]
It forces $p\ge19$. Both signs therefore survive the peels, and an
opposite pair is present.

Put
\begin{gather*}
 d=|D|=92,\quad e=|E|=84,\quad E=P\cup\{0\}\cup(-Q),\quad
 a=|P|,\ b=|Q|,\ a+b=83,\\
 \{a,b\}=\{41-k,42+k\},\quad I=k(k+1),\quad
 \Lambda=1+T_a+T_b=1765+I,\quad \Omega=356-I.
\end{gather*}
For a positive set $X$ of size $r$ let
\[
 \mu(X)=|\FS(X)|-1-T_r,\quad
 \delta_0(X)=|\{x\in X:2x\notin R(X\cup\{0\})\}|,
 \quad\rho(X)=\mu(X)+2-2\delta_0(X),
\]
where $R$ denotes sums of two distinct elements. Set
$\rho=\rho(P)+\rho(Q)$, and let
\[
 c=2|2D|+18d+85-|\FS(D)|,\quad
 \kappa=|\FS(E)|-\Lambda-\mu(P)-\mu(Q).
\]
The inherited nonnegative slacks satisfy
\begin{equation}\label{eq:budget}
 \rho+2\theta+c+\T=356-I.
\end{equation}
Here $\T=\kappa+\ell_1+\ell_2+g$, with the same peeling slacks as in
\cite{base100}; $\theta\ge0$ is the diagonal-exception slack.

For every oriented $a,b\ge19$ with $a+b=83$ and $\Omega\ge0$, check
\[
 T_{\min(a,b)}-2-\Omega>0,\qquad
 2(a+b)+4e+14-\Lambda+3\Omega<0,
 \qquad \Lambda-(4d-10+\Omega)>0.
\]
There are 38 orientations. Their first and third minimum margins are
230 and 1051; the second expression has maximum $-181$.
Thus the inverse theorem and propagation argument in \cite{base103}
apply at this size. After common-lattice scaling, $D$ is integral.
No integrality of the eight outer elements is asserted at this stage.

Let $h_P,h_Q,z,w$ be the hull-hole counts of the supports of $P,Q,E,D$
on the common lattice. The same structural argument yields
\begin{gather}
 h_P,h_Q\le\kappa,\quad z=h_P+h_Q-\kappa,
 \quad0\le z\le\min(h_P,h_Q)\le\kappa,
 \quad h_X<|\FS(X)|-1,\label{eq:holes}\\
 w\le4z,\qquad
 \T\ge\kappa+172-3z,\quad
 \T\ge\max\{58,z,172-2z\}.\label{eq:T}
\end{gather}
All side sizes exceed the inherited thresholds for the estimates
$\rho(X)\ge0$ and the span lemma. If $z>0$, neither side support is a
progression, so $\mu(P)\ge a-3$, $\mu(Q)\ge b-3$, and $\rho\ge55$.

\section{High holes and the all-active cases}
We reuse the parameterized effective-hole lemma proved in \cite{base100}.
For a positive set of size $r$, a valid witness $(r,C,\beta)$ proves
$\mu\ge r+C+1$ whenever $2\beta<h<|\FS|-1$.
For completeness its exact arithmetic tests are
\begin{gather*}
 K=\max\{k\ge3:k(k-3)/2\le C+1\},\quad D_*=C+K,\quad E_*=D_*+1,\\
 M=\max\left\{2\beta,1+\left\lfloor
 \frac{1+\sqrt{9+8(E_*+6\beta)}}2\right\rfloor\right\},\quad
 J=\max\{K+1,2\beta+5,M\},\\
 T_M-6\beta>M+1+E_*,\qquad
 \beta(r-J+1)>D_*,\qquad T_{J-1}-1>D_*.
\end{gather*}
As before, search only $20\le r\le64$, $7\le\beta\le20$, $0\le C\le600$.
Each chosen witness is rechecked exactly; no assertion of optimality of
this finite witness search is needed.

For every integer $15\le z\le\Omega$, use the strongest available row
with $2\beta<z$ for each side. If its constant is $C$, put
\[
 u_r=r+C+1,\qquad v_r=\max\{r-14,u_r+2-2r\}.
\]
The bounds $h_P,h_Q\ge z$ license the lemma. The second inequality uses
$\delta_0(X)\le r$. Every actual high-hole state satisfies
\begin{equation}\label{eq:cmax}
 c\le C_k(z):=356-I-v_a-v_b-\max\{58,z,172-2z\}.
\end{equation}
Negative values discard a state. The upper limit $z\le\Omega$ follows
from $z\le\kappa\le\T\le\Omega$, so the loop is exhaustive.

Write $R_p=\max P$, $R_q=\max Q$, $M_p=\max D$, $M_q=-\min D$,
and let $S_p,S_q$ be the positive and negative-magnitude masses in $D$.
Set
\[
 \varepsilon_p=\sum P-T_a=\mu(P)+h_P,\quad
 \varepsilon_q=\sum Q-T_b=\mu(Q)+h_Q,
\]
\[
 F=S_p+S_q+c-2(M_p+M_q)-2(R_p+R_q).
\]
The inherited covering estimates give
\begin{equation}\label{eq:Fupper}
 F\le18d+126+w\le1782+4z.
\end{equation}
If all eighteen ordinary covering units are nonempty, after a possible
reflection one has
\begin{align}
 F&\ge F_A(a,b)+\tfrac12\varepsilon_p+\varepsilon_q,
 &F_A&=\tfrac12T_a+T_b+6a+8b+85,\label{eq:FA}\\
 F&\ge F_B(a,b)+\tfrac23\varepsilon_p+\varepsilon_q,
 &F_B&=\tfrac23T_a+T_b+6a+\tfrac{22b+236}{3}.\label{eq:FB}
\end{align}
The theorem is applicable to any all-active configuration; it does not
require $c<d$ if activity is known directly.

For high holes use $\varepsilon_p\ge u_a+z$, $\varepsilon_q\ge u_b+z$.
Every surviving state has $C_k(z)\le64<92$, which forces activity of all
units. Both lower bounds exceed \eqref{eq:Fupper}.
For low holes use $u_a=u_b=0$ at $z=0$ and
$u_a=a-3$, $u_b=b-3$ otherwise, and the corresponding lower bounds
$\rho\ge0$ or $\rho\ge55$. The exact summary is
\begin{center}
\begin{tabular}{lrrr}
\toprule
Core holes&States&Minimum active margin&High-hole $c$ upper bound\\
\midrule
$0\le z\le14$&364&33&--\\
$z\ge15$&966&150&64\\
\bottomrule
\end{tabular}
\end{center}
The minima include both orientations and both alternatives. Consequently
all high-hole states and every low-hole state with no empty unit are
excluded. Of the 364 low-hole states, 222 allow an empty unit; the next
sections cover all their configurations. Scalar states are necessary
relaxations, not asserted realizable sets.

\section{Empty units, outside centers, and equality}
At zero holes, the stronger first inequality in \eqref{eq:T} gives
\begin{equation}\label{eq:zero}
 \rho+\kappa+2\theta+c\le184-I.
\end{equation}
At positive low holes,
\begin{equation}\label{eq:positive}
 c\le129-I+2z\le157-I.
\end{equation}
An empty unit costs $d=92$ in the covering slack. Thus at most two units
are empty, and two are possible only at zero holes.

\subsection{The ordinary centers and exhaustive modes}
Write the outer magnitudes as $p_1>p_2>p_3>p_4>M_p$ and
$q_1>q_2>q_3>q_4>M_q$. Define
\[
 \lambda_p=p_1-p_2-p_3-p_4,\quad
 \alpha_p=p_2-p_3-p_4,\quad
 \beta_p=2p_2-p_1-p_3-p_4,
\]
and reflect to define the $q$ quantities. Index the ten positive centers
by $0,\ldots,9$:
\[
 v^p=(\lambda_p+p_1,\ldots,\lambda_p+p_4,
       \lambda_p-q_1,\ldots,\lambda_p-q_4,\alpha_p,\beta_p).
\]
Interchange $p,q$ for the negative centers $v^q$; a negative center
translates $-D$ into the target $[1,S_q]$.
Slot labels $p_i,q_i$ in the mode lists refer to these zero-based centers,
and are distinct from the four magnitude variables with subscripts 1--4.

The first paired unit uses centers $p4,q4$. In branch A,
$\beta_p+\beta_q\le0$, the second uses $p9,q9$; in branch B,
$\beta_p+\beta_q>0$, it uses $p5,q5$. The other sixteen units are singles.
Choose the missing unit or unordered pair of missing units. Retain all
other singles. If the first pair is present, reflect to choose its
positive contributor $p4$. If the second pair is present, retain its
positive and negative contributor in separate modes, except in branch A
when either single $p5$ or $q5$ is missing: in that situation drop all
individual constraints on the second pair and do not duplicate the mode.
Dropping the pair allows it to split and enlarges the model.

These choices are exhaustive. The first pair always has at most one
contributor because its center sum is less than $-(M_p+M_q)$. The same
holds for the second pair in B, and in A whenever $p5,q5$ both contribute,
using
\[
 \beta_p+\beta_q=-(v^p_5+v^q_5)-2(p_3+p_4+q_3+q_4).
\]
There are 33 A modes and 35 B modes with one missing unit; with two
missing units there are 258 A modes and 289 B modes. If the first pair
is missing both second-pair choices are still retained.

\subsection{One empty unit}
Put $\delta=c-92$. The outer-integrality argument of \cite{base100}
applies with exactly one empty unit: all outer magnitudes are integral,
so their strict gaps are at least one. For a retained positive center
let $l_v,u_v$ count points of $v+D$ below 1 and above $S_p$. Then
\begin{equation}\label{eq:loss}
 \sum_{v\text{ retained}}(l_v+u_v)\le\delta.
\end{equation}
For an in-hull center,
\begin{equation}\label{eq:quantile}
 -v-l_v\le-b-5,\qquad v-u_v-S_p\le-a-4.
\end{equation}
The valid replacements at an outside endpoint are
\begin{align}
 v<0:&\quad -v-l_v-M_p\le-92,\quad l_v\ge b+5,\label{eq:lower-out}\\
 v>S_p:&\quad v-u_v-S_p-M_q\le-91,\quad u_v\ge a+5.\label{eq:upper-out}
\end{align}
Reflect these inequalities for negative centers. The opposite ordinary
endpoint row remains valid. In particular each outside center costs at
least $L=\min(a,b)+5=46-k$.

At zero holes $\delta\le92-I$; at positive low holes
$\delta\le65-I$. Their nonnegativity gives $k\le9$ and $k\le7$,
respectively. One outside center requires $L\le\delta$, giving
$k^2\le46$ or $k^2\le19$. Moreover,
\[
 2L-(65-I)=(k-\tfrac12)^2+\tfrac{107}{4}>0,
 \quad 3L-(92-I)=(k-1)^2+45>0.
\]
There is at most one outside center with positive low holes, and at most
two at zero holes. The all-in-hull regimes are Z and P, the one-outside
regimes ZE and PE, and the two-outside regime is ZZ.

\subsection{Rigidity in the two-outside case}
In ZZ, $2L\le\delta\le92-I$ forces $k(k-1)\le0$. Thus $k=0$ or 1,
$\delta=2L$, and \eqref{eq:zero} forces
\begin{equation}\label{eq:rigid}
 \rho(P)=\rho(Q)=\kappa=0.
\end{equation}
Both core supports are full integer intervals by \eqref{eq:holes}.

\begin{lemma}[Zero-defect core]\label{lem:core}
Let $X\subset\Z_{>0}$ have $r\ge15$ elements. If
$\FS(X)=[0,\sum X]_{\Z}$ and $\rho(X)=0$, then
\[
 X=[1,r]_{\Z}\quad\hbox{or}\quad
 X=[1,r+1]_{\Z}\setminus\{r-1\}.
\]
In particular $\max X\le r+1$.
\end{lemma}
\begin{proof}
The inherited prefix recurrence
$\rho(X_r)\ge\rho(X_{r-1})+y_r-2$ and
$\rho(X_{r-1})\ge0$ give $y_r\le2$.
The small-increment inverse lemma therefore writes $X$, up to scaling,
as $[1,r+q]_{\Z}\setminus D_0$ with $0\le q\le2$ deletions and maximum
$h=r+q$ present. This deleted interval is primitive since it has more
than half the points of its hull. Since $X$ is integral, its scale is an
integer; full support includes 1, so that scale equals one.

Let $C=\{h-d:d\in D_0\}$. Full support and $\rho=0$ imply
\[
 \mu=\sum C-T_{q-1}=2\delta_0(X)-2.
\]
Here the $q=0$ case is immediate and is treated separately, so the
display is used for $q\ge1$. The deleted-interval defect bound gives
$\delta_0(X)\le2q$, hence $\mu\le6$ for $q=1,2$.
A low defect $x\le q$ forces the hole $2x$ and a depth at least $r-2$
in $C$, giving $\mu\ge r-3>6$, a contradiction. There are no low defects.
All other defects lie at the upper end. Besides the maximum $h$, let
them be $h-t$ for $m$ distinct values of $t$, where $m\le q$.
Presence and defectivity give $t\notin C$, $2t\in C$.
Replacing these distinct depths $2t$ by $t$ leaves $q$ distinct positive
integers. Consequently
\[
 2m=\mu\ge q+\sum t\ge q+T_m,\qquad 0\le m\le q\le2.
\]
The only nontrivial possibility is $q=m=1$. It has $\mu=2$ and
$C=\{2\}$, hence $D_0=\{r-1\}$. This proves the claim.
\end{proof}

By the lemma, $R_p-a,R_q-b\le1$ in ZZ. Each retained center loses at
most $L\le46$ points: another outside center already costs $L$, or both
outside centers are distinct from it and exhaust the budget. Its hull
count is therefore at least $92-46=46>4$.
If a positive center is below zero, at least one of these contributing
points comes from the positive core, since only four positive points
of $D$ are outside $P$. Thus $-v<R_p$. Counting missing integer
positions in $[1,R_p]$, at most $R_p-a\le1$, gives
$-v-l_v\le-b-4$. If $v>S_p$, the reflected core argument gives
$v-u_v-S_p\le-a-3$. In-hull centers already obey stronger rows.
Thus every retained positive center simultaneously satisfies the relaxed
bounds
\begin{equation}\label{eq:ZZ}
 -v-l_v\le-b-4,\qquad v-u_v-S_p\le-a-3,
 \quad R_p\le a+1,\quad R_q\le b+1.
\end{equation}
Their reflections cover the negative side. This avoids enumerating the
locations of the two exceptions.

\subsection{Two empty units: full translates without outer integrality}
Their cost is 184. Equations \eqref{eq:zero}--\eqref{eq:positive} force
$z=0$, $k=0$, $\{a,b\}=\{41,42\}$, $c=184$ and all residual cover
losses zero. Every remaining unit reaches its full capacity of 92.
Each retained individual translate therefore lies wholly in its actual
target support, hence in its target interval. Such a center is integral
because it has a contribution between integral points. For a positive
center this yields
\begin{equation}\label{eq:full}
 v\ge M_q+1,\qquad v\le S_p-M_p,
\end{equation}
with reflected inequalities on the negative side.

This does not require proving the eight outer magnitudes integral.
In the resulting EE models we use only
$p_i\ge p_{i+1}$, $q_i\ge q_{i+1}$, $p_4\ge M_p$, $q_4\ge M_q$,
and replace the strict B condition by $\beta_p+\beta_q\ge0$.
The integral inner set still gives $M_p\ge R_p+4$ and its reflection.
Retained pair contributors and possible split pairs are handled by the
exhaustive rule above. There are $547\cdot2=1094$ EE systems.

\section{The six complete finite model families}
The table specifies all integer side sizes, budgets, and special rows.
Both orientations are included by taking every integer $a$ in the stated
interval and $b=83-a$.
\begin{center}\small
\begin{tabular}{lrrrrl}
\toprule
Regime&Range of $a$&$\delta$ upper bound&Upper $F$&Systems&Extra information\\
\midrule
Z&32--51&$92-I$&1782&1360&All retained centers in-hull\\
P&34--49&$65-I$&1838&1088&In-hull; mass bonuses $a-2,b-2$\\
ZE&35--48&$92-I$&1782&32312&One designated outside endpoint\\
PE&37--46&$65-I$&1838&23080&One outside endpoint; mass bonuses\\
ZZ&40--43&$92-I$&1782&272&Core bounds and \eqref{eq:ZZ}\\
EE&41--42&0&1782&1094&Two missing units; full translates\\
\bottomrule
\end{tabular}
\end{center}
In EE, $\delta$ denotes the residual slack $c-184$, which is zero.
In the other regimes it is $c-92$.
ZE and PE choose every retained center and each of its two endpoints.
Sixty-six one-missing-unit modes retain 17 centers and two retain 16,
giving $2(66\cdot17+2\cdot16)=2308$ endpoint choices for each $a$.
All six families together contain 59,206 systems.

\subsection{Every row of the systems}
All variables
\[
 x=(p_1,\ldots,p_4,q_1,\ldots,q_4,M_p,M_q,R_p,R_q,S_p,S_q,
       \delta,(l_v,u_v)_{v\text{ retained}})
\]
are nonnegative real numbers. The model relaxes the original integer
variables except the explicitly enumerated parameter $a$.
Rows are oriented $Ax\le b$. Include the reflected counterpart of each
positive-side row:
\begin{center}\small
\begin{tabular}{ll}
\toprule
Row&Constraint\\
\midrule
$D_i^p$ ($1\le i\le3$)&$p_{i+1}-p_i\le-1$; right side 0 in EE\\
$D_4^p$&$M_p-p_4\le-1$; right side 0 in EE\\
$E_p,G_p$&$R_p-M_p\le-4$, $-R_p\le-a$\\
$H_p$&$2M_p+6R_p-2S_p\le-2T_a-12-2\epsilon_p$\\
cap&$\delta$ at most the tabled upper bound\\
branch A&$\beta_p+\beta_q\le0$\\
branch B&$-\beta_p-\beta_q\le-1$; right side 0 in EE\\
loss&$\sum(l_v+u_v)-\delta\le0$\\
Fupper&$S_p+S_q+\delta-2(M_p+M_q+R_p+R_q)\le F_{\max}-92j$\\
\bottomrule
\end{tabular}
\end{center}
Here $j=2$ in EE and $j=1$ otherwise. The mass bonuses are
$\epsilon_p=a-2$, $\epsilon_q=b-2$ in P and PE, and zero otherwise.
The mass rows follow from the exact inequalities
\[
 S_p\ge T_a+\varepsilon_p+3R_p+6+M_p,\qquad
 S_q\ge T_b+\varepsilon_q+3R_q+6+M_q.
\]
They retain the exact triangular numbers and the same imbalance $I$ as
the budget; neither is replaced by a global minimum.

For ordinary retained centers use \eqref{eq:quantile}; in ZE and PE
replace the designated endpoint row by \eqref{eq:lower-out} or
\eqref{eq:upper-out} and add its unavoidable-loss row. Its opposite row
is unchanged. The actual outside-location inequality is omitted, which
only weakens the system. In ZZ use \eqref{eq:ZZ}.
In EE the rows are
\[
 -v-l_v+M_q\le-1,\qquad v-u_v-S_p+M_p\le0,
\]
and their reflections. Since the loss row and cap give all losses zero,
these are exactly \eqref{eq:full}. No strict unit outer gap is smuggled
into this regime.

\subsection{Integer certificates and complete interval coverage}
For each system the supplied proof data give a nonnegative integer row
vector $w$ such that
\begin{equation}\label{eq:farkas}
 wA\ge0\quad\text{coefficientwise},\qquad wb<0.
\end{equation}
This proves inconsistency: $0\le wAx\le wb<0$ for any $x\ge0$.
No numerical tolerance is used.

There are 204 distinct vectors in
\path{data/upper99_certificates.json}. The model generator produces 5,367
structural groups, each consisting of a regime, a missing-unit mode and,
when needed, an outside endpoint. For a group the proof data partition
its entire integer $a$ interval into consecutive subintervals assigned
to vectors. This is a storage compression only: the verifier rebuilds
the system and checks \eqref{eq:farkas} for every integer $a$ individually.
It requires equality of the generated and stored group-key sets, then
rejects gaps, overlaps, extra parameters and missing endpoints in every
partition. It also rejects weights on absent rows.

All 59,206 systems pass. The least normalized contradiction gap
$-wb/\sum_iw_i$ is $12/83>0$. The report's table and matrix recipe, rather
than a stored list of optimizer statuses, define the model universe.
The optional discovery script uses SciPy to propose weights, converts
them to rational numbers and then integers, and checks their coefficients
before saving. It also invokes the full coverage checker before writing
the output. The delivered verification path does not import SciPy or
call an optimizer.

\section{Closing the bound and reproducibility}
\begin{proof}[Proof of Theorem~\ref{thm:main}]
The previous theorem gives $N\le100$. If a 100-element feasible set
existed, adjoining-zero and integralization arguments would put it in
the setup of Section 2. The strict gates establish the common lattice.
Section 3 excludes high holes and all low-hole configurations with no
empty unit. Equations \eqref{eq:zero}--\eqref{eq:positive} leave at most
two empty units. With one empty unit, the outside-center count partitions
the configuration into Z, P, ZE, PE or ZZ; the equality lemma justifies
the stronger ZZ rows. With two empty units, full translates and weak
outer order put it in EE. Every corresponding system is inconsistent
by \eqref{eq:farkas}. Thus size 100 is impossible and $N\le99$.
No step assumes feasibility passes to arbitrary subsets.
\end{proof}

Run from the repository root with Python 3.10 or later:
\begin{verbatim}
python -I -S -B verification/upper100.py --output .work/base100.json
python -I -S -B verification/upper99.py --output .work/upper99.json
\end{verbatim}
The new program uses the standard-library arithmetic helpers in
\path{verification/upper100.py}. It reads only its new certificate data,
not archived endpoint data. The report, complete run and input hashes
are recorded under \path{results/2026-10-16/upper99/}.
An optimized-mode run with \texttt{-O} gives identical JSON; all checks
use explicit exceptions rather than removable assertions.

Seven negative controls reject an invalid hole row, negative weights,
a noncontradictory positive row, an unknown row, a missing group, a missing
integer inside a coverage interval, and a request to certify size 99
using the size-100 routine. In addition, a separate exact enumeration
checks 4,776 deleted intervals with $15\le r\le30$ and at most two
deletions. Its 32 zero-defect full-support cases are precisely the two
families in Lemma~\ref{lem:core} for every tested $r$. This is finite
corroboration; the infinite equality statement is proved in the lemma.

Optional fresh certificate discovery, requiring SciPy, is
\begin{verbatim}
python experiments/generate_upper99_certificates.py \
  --output .work/regenerated_upper99.json
python -I -S -B verification/upper99.py \
  --certificates .work/regenerated_upper99.json
\end{verbatim}
The PASS result certifies the displayed finite reductions. Applicability
to real sets relies on the written structural arguments and their
explicitly inherited foundations, not on the finite run alone.

At the next size, $n=99$, the balanced zero-hole parameters are
\[
 a=b=41,\quad d=91,\quad\Omega=376,\quad G=170,
 \quad c\le206.
\]
Two empty units cost 182 and can leave residual slack 24. The zero-loss
argument used for EE is consequently unavailable without further work.
This is the next obstruction to this particular proof, not a claimed
99-element feasible set.

\begin{thebibliography}{9}
\bibitem{base103} Sum-of-Subset project,
\emph{A checked continuation from the upper bound 115 to 103},
14 October 2026. \path{reports/upper103_from115_20261014.tex}.
\bibitem{base100} Sum-of-Subset project,
\emph{A checked continuation from 103 to 100}, 16 October 2026.
\path{reports/upper100_from103_20261016.tex}.
\bibitem{archive} Sum-of-Subset project,
\emph{Complete upper-86 report}, 5 October 2026.
\path{reports/upper86_complete_20261005.tex}, section ``Excluding size 100''
and its upper-99 certificate appendix. Source of the endpoint strategy;
the full archived chain is not certified by this continuation.
\end{thebibliography}
\end{document}
